% TOP_PROBLEM: {"id": 1286, "title": "Beal's Conjecture", "category_id": 1, "category": {"id": 1, "name": "number_theory", "display_name": "Number Theory", "description": "Properties of integers, prime numbers, Diophantine equations.", "slug": "number-theory", "order_index": 1, "created_at": "2026-07-31T15:26:25.670Z"}, "status": "open"}
% FIRST_POSTED: 2026-09-27
% !TeX program = lualatex
% ATTEMPT_STATUS: unresolved
\documentclass[11pt]{article}
\usepackage[margin=25mm]{geometry}
\usepackage{fontspec}
\setmainfont{Latin Modern Roman}
\usepackage{amsmath,amssymb,mathtools}
\usepackage{unicode-math}
\setmathfont{Latin Modern Math}
\usepackage{xurl,hyperref}
\hypersetup{colorlinks=true,urlcolor=blue}
\setcounter{secnumdepth}{0}
\setlength{\parindent}{0pt}
\setlength{\parskip}{5pt}
\setlength{\emergencystretch}{3em}
\begin{document}
\section*{\texorpdfstring{Beal's Conjecture}{Beal's Conjecture}}\label{1286}
\subsection{Definitions and mathematical statement}
Beal's conjecture asserts that for positive integers $A,B,C$ and exponents $x,y,z>2$,
\[
 A^x+B^y=C^z\quad\Longrightarrow\quad\gcd(A,B,C)>1.
\]
Equivalently, some prime must divide all three bases. If a prime divides two bases in such an equation, it divides the third, so excluding a common prime also enforces pairwise coprimality. The exponents may be different and need not be prime. Positivity excludes zero and signed degeneracies. Unlike a finiteness statement for exceptional solutions, the assertion prohibits every coprime solution in this exponent range.
\par\smallskip\textit{Formulation and concept references:} \href{https://sites.math.unt.edu/~mauldin/beal.html}{[S1]}.

\subsection{Short English statement}
Whenever the sum of two positive perfect powers equals a third and all exponents exceed two, must their bases share a prime factor?
\subsection{Research attempt}
\textbf{Status: UNRESOLVED.} No primitive solution is constructed and no proof excluding every primitive solution is obtained. The attempt reduces the possible exponent signatures, gives a complete elementary descent for a substantial quartic subcase, and proves a conditional finiteness statement from the $abc$ conjecture while explaining why finiteness falls short of Beal's assertion.

\paragraph{Statement and source audit.}
The equation in the catalogue is the positive-integer equation
\[
 A^x+B^y=C^z,\qquad x,y,z>2.
\]
The formulation on Mauldin's page [S1], checked on 27 September 2026, asks that the bases have a common prime divisor. It allows unequal exponents and bases equal to one. Throughout this attempt a primitive solution means $\gcd(A,B,C)=1$. Results imported from Fermat's Last Theorem and Mih\u{a}ilescu's theorem are explicitly identified below; their deep proofs are not replaced by elementary manipulations. Any use of $abc$ is conditional.

\paragraph{1. Primitivity, signs, and the unit bases.}
If a prime divides any two of $A,B,C$, the equation forces it to divide the third. For example, $q\mid A,B$ gives $q\mid C^z$ and hence $q\mid C$. The other two cases follow by rearranging. Therefore a primitive solution is automatically pairwise coprime. Conversely, pairwise coprimality implies primitivity. This equivalence depends on the equation; it is not true for arbitrary triples of integers.

The case $C=1$ is impossible because the left side is at least two. If $A=B=1$, the equality $C^z=2$ is also impossible. If exactly one of $A,B$ is one, the equation becomes a difference of two perfect powers equal to one. The classical theorem of Mih\u{a}ilescu says that the only consecutive nontrivial perfect powers, with both exponents greater than one, are $8$ and $9$. Their upper exponent is two, so this cannot occur with both $y,z>2$ (or both $x,z>2$ after interchanging the summands). Thus, using that external theorem, every hypothetical primitive solution has $A,B,C\ge2$.

Positivity matters in the estimates later: both $A^x$ and $B^y$ are strictly smaller than $C^z$. An argument allowing arbitrary signed bases or moving a negative term without checking parity would be treating a different equation.

\paragraph{2. Reduce exponents individually, not to a common exponent.}
Every integer $m>2$ either has an odd prime divisor or is divisible by four. Choose a divisor $r$ of $x$, a divisor $s$ of $y$, and a divisor $t$ of $z$, each belonging to
\[
 \{4\}\cup\{\text{odd primes}\}.
\]
Then
\[
 (A^{x/r})^r+(B^{y/s})^s=(C^{z/t})^t.
\]
Raising a positive integer to a positive power does not change its prime support, so pairwise coprimality and primitivity are preserved. Hence excluding primitive solutions for every signature $(r,s,t)$ of this restricted type would prove the conjecture.

This reduction still leaves infinitely many mixed signatures. It is invalid to replace the three chosen divisors by a single common odd prime when no such divisor exists. In particular, signatures such as $(3,4,5)$ survive untouched. Fermat's Last Theorem concerns equal exponents; it does not settle all equations whose three exponents happen to be at least three.

There is a correct common-divisor subcase. If $d=\gcd(x,y,z)>2$, then
\[
 (A^{x/d})^d+(B^{y/d})^d=(C^{z/d})^d
\]
contradicts Fermat's Last Theorem. Thus there are no positive solutions at all in that subcase. If $d=2$, the transformed equation is a Pythagorean equation, for which positive solutions exist, so the same inference is unavailable. If $d=1$, it gives no restriction.

A common factor of the bases also cannot simply be divided out while preserving the original exponents. Writing $A=aD,B=bD,C=cD$ introduces different powers $D^x,D^y,D^z$. Conversely, if any solution exists, then for $L$ a common multiple of $x,y,z$ and an integer $u>1$,
\[
 (u^{L/x}A)^x+(u^{L/y}B)^y=(u^{L/z}C)^z
\]
is another solution with a common base divisor. This explains why producing large nonprimitive solutions is easy and irrelevant to disproving the assertion. For example, $2^3+2^3=2^4$ and $3^3+6^3=3^5$ satisfy the equation and its predicted common-divisor conclusion.

\paragraph{3. Complete descent for $X^4+Y^4=Z^2$.}
We prove that there are no positive integers satisfying this auxiliary equation. If a solution exists, divide $X,Y$ by their greatest common divisor $d$. Since $d^4\mid Z^2$, valuation comparison gives $d^2\mid Z$, so this produces a primitive solution. Choose a primitive solution with $Z$ minimal. The two bases cannot both be odd, because fourth powers of odd integers are one modulo 16 and a square cannot be two modulo 16. They cannot both be even by primitivity. Interchanging them, take $X$ odd and $Y$ even.

The primitive Pythagorean triple $(X^2,Y^2,Z)$ therefore has the form
\[
 X^2=m^2-n^2,\qquad Y^2=2mn,\qquad Z=m^2+n^2,
\]
with $m>n>0$, $\gcd(m,n)=1$, and opposite parity. As $X^2\equiv1\pmod4$, the first identity forces $m$ odd and $n$ even. Unique factorization in $2mn=Y^2$ now gives
\[
 m=a^2,\qquad n=2b^2,\qquad Y=2ab,
 \qquad \gcd(a,b)=1.
\]
Indeed every odd prime appears with even valuation in each of the coprime factors, while the 2-adic valuation of $n$ is odd. Substitution yields
\[
 X^2+(2b^2)^2=a^4.
\]
This is another primitive Pythagorean triple: $X$ is odd, and a prime dividing $X$ and $b$ would divide $a$ from the displayed identity. Parameterize it as
\[
 X=r^2-s^2,\qquad 2b^2=2rs,\qquad a^2=r^2+s^2,
\]
with $r>s>0$ and $\gcd(r,s)=1$. Since $b^2=rs$ and the factors are coprime, $r=c^2$ and $s=d^2$ for positive integers $c,d$. Thus
\[
 c^4+d^4=a^2
\]
is a new solution of the original auxiliary equation, with smaller last variable: $a<Z=a^4+4b^4$. This contradicts minimality. The descent is complete.

It follows that the Beal equation has no positive solutions whenever
\[
 4\mid x,\qquad 4\mid y,\qquad 2\mid z.
\]
Set $X=A^{x/4},Y=B^{y/4},Z=C^{z/2}$ and apply the proved result. This includes some mixed signatures which the common-exponent argument alone misses. It does not cover, for example, a cubic summand with an odd exponent on the other summand. The descent relies twice on one leg already being a square and on a particular even leg having the form $2b^2$; those identities cannot be assumed for an arbitrary mixed signature.

\paragraph{4. What $abc$ would prove uniformly.}
Assume the following standard $abc$ statement: for every $\delta>0$ there is a constant $K_\delta$ such that pairwise coprime positive integers $a+b=c$ satisfy
\[
 c\le K_\delta\operatorname{rad}(abc)^{1+\delta},
\]
where $\operatorname{rad}$ denotes the product of distinct prime divisors. Apply this to a primitive Beal solution, with $a=A^x,b=B^y,c=C^z=M$. Its radical is $\operatorname{rad}(ABC)$, bounded above by $ABC$. Since $A<M^{1/x}$, $B<M^{1/y}$, and $C=M^{1/z}$,
\[
 M\le K_\delta M^{(1+\delta)\sigma},\qquad
 \sigma=\frac1x+\frac1y+\frac1z.
\]
The signature $(3,3,3)$ has no solution by Fermat's Last Theorem. In every other signature with exponents at least three, at least one exponent is at least four, so
\[
 \sigma\le\frac13+\frac13+\frac14=\frac{11}{12}.
\]
Choose the fixed number $\delta=1/22$. Then $(1+\delta)\sigma\le23/24$, and consequently
\[
 M^{1/24}\le K_{1/22},\qquad M\le K_{1/22}^{24}.
\]
This bound is uniform over all remaining exponent signatures. After the unit bases have been excluded, $A,B,C\ge2$, so each exponent is at most $\log_2M$ and all bases are bounded by $M$. Hence $abc$, together with the named established theorems, implies that only finitely many primitive Beal solutions can exist.

The conclusion is finiteness, not emptiness. An unknown $abc$ constant does not give a completed finite search, and even an effective bound would require checking every case below it. It is therefore incorrect to announce Beal's conjecture proved merely from this exponent-sum argument. There is also no unconditional $abc$ input being asserted here: the whole paragraph is a conditional implication, with the exact assumption and its use displayed.

\paragraph{5. Congruences and finite searches.}
Congruences can discard base residue patterns for a fixed signature, but a locally admissible residue pattern need not lift to a positive integer solution. The obvious residues $A\equiv0,B\equiv C\equiv1$ modulo an auxiliary modulus already show why some patterns survive many simple tests. They do not satisfy a common-prime conclusion as an integer identity, and local admissibility supplies no global counterexample.

For reproducibility, the offline diagnostic uses exact integer powers, not floating-point roots. It stores every value $a^e$ with $1\le a\le100$ and $3\le e\le8$, retaining all representations of values such as $64$. It adds each pair of stored values and looks up every representation of the sum in the same table. It checks the common-divisor conclusion for each resulting equation, includes the displayed nonprimitive examples, and records the exact finite box searched. A separate check enumerates primitive Pythagorean parameters and verifies the algebra and parity identities used in the descent. No result outside that box is inferred from the computation.

Perfect-power representations must not be collapsed to a single arbitrary exponent: $a^e=b^f$ can lead to different exponent signatures, and the problem quantifies over all of them. Likewise, the search cutoff on $C$ is a cutoff on a base, not on the entire integer $C^z$. These distinctions are included to keep the computational claim checkable.

\paragraph{6. The remaining gap.}
The proven reductions leave primitive mixed signatures with no common exponent divisor greater than two and outside the quartic descent pattern. Addressing all of them would require a uniform obstruction, or a reduction to finitely many cases with a fully effective and completed verification. Neither is provided. The argument does establish exactly why the equal-exponent theorem, quartic descent, and conditional radical bound are relevant, and exactly why their union is still insufficient. The unrestricted assertion remains unresolved by this attempt.

\subsection{Sources}
\begin{itemize}
\item[S1]R. D. Mauldin, The Beal Conjecture and Prize, University of North Texas; announced in Notices Amer. Math. Soc. 44 (1997), 1436-1437.
\url{https://sites.math.unt.edu/~mauldin/beal.html}.
\textit{Formulation source.}
\item[E1]Preda Mih\u{a}ilescu, Primary cyclotomic units and a proof of Catalan's conjecture, Journal f\"ur die reine und angewandte Mathematik 572 (2004), 167--195.
\url{https://doi.org/10.1515/crll.2004.048}.
\textit{Added primary source: The theorem on consecutive perfect powers, used only for unit bases. Publisher bibliographic record consulted 27 September 2026.}
\item[E2]Andrew Wiles, Modular elliptic curves and Fermat's Last Theorem, Annals of Mathematics 141 (1995), 443--551.
\url{https://annals.math.princeton.edu/1995/141-3/p01}.
\textit{Added primary source: Established equal-exponent theorem, used as external input rather than reproved. Consulted 27 September 2026.}
\end{itemize}
\end{document}
